% ===== FILE 1/9: Bìa, Động lực, 3D Gaussian, Hiệp phương sai, Scale khởi tạo, Scene extent, Màu SH, Lịch tăng bậc SH =====

% --- Slide: Bìa ---
\begin{frame}
\titlepage
\end{frame}

% --- Slide: Động lực ---
\begin{frame}{Vì sao 3D Gaussian Splatting?}
\footnotesize
\begin{itemize}
    \item \textbf{Bài toán đặt ra:} cho một tập ảnh chụp một cảnh từ nhiều góc, cùng camera pose ước lượng sẵn bởi
    SfM/COLMAP, hãy dựng lại một biểu diễn 3D sao cho có thể \textbf{render lại ảnh từ bất kỳ góc nhìn mới nào}
    (novel view synthesis) một cách chân thực.
    \item \textbf{NeRF (2020)} giải bài này bằng một mạng MLP ẩn ánh xạ (vị trí, hướng nhìn) $\to$ (màu, mật độ), rồi
    \textbf{ray-marching}: bắn hàng trăm điểm mẫu dọc mỗi tia sáng, gọi mạng hàng trăm lần cho mỗi pixel. Chất lượng
    cao nhưng cực chậm --- train mất hàng giờ/ngày, render cũng mất giây/khung hình vì phải lặp lại forward pass qua mạng cho từng điểm mẫu.
    \item \textbf{3D Gaussian Splatting} (Kerbl, Kopanas, Leimkühler, Drettakis --- SIGGRAPH 2023) đổi hẳn cách biểu diễn:
    thay mạng ẩn bằng \textbf{hàng trăm nghìn tới hàng triệu "Gaussian 3D" tường minh}, mỗi Gaussian là một đối tượng
    có tham số cụ thể $\{\mu,\Sigma,\alpha,c\}$ (Slide sau). Render không cần mạng nơ-ron, không cần ray-marching --- chỉ
    cần \textbf{chiếu (project) và tô (rasterize)} trực tiếp các Gaussian lên màn hình, giống hệt cách GPU vẽ tam giác
    trong đồ hoạ game truyền thống. Kết quả: train nhanh hơn nhiều lần, và quan trọng nhất --- render \textbf{thời gian thực}
    (hàng chục--hàng trăm khung hình/giây) sau khi train xong.
    \item \textbf{Nhưng còn một vấn đề chưa giải quyết trọn vẹn:} số lượng Gaussian $N$ cần bao nhiêu, đặt ở đâu, kích
    thước ra sao --- không biết trước, phải \textbf{tự động điều chỉnh trong lúc train} (đây gọi là \emph{Adaptive Density
    Control}, viết tắt ADC). Cơ chế ADC gốc chỉ dựa vào \textbf{gradient} để quyết định thêm/bớt Gaussian --- và đây
    chính là chỗ \textbf{SADGS} (Structure-Aware Densification for Gaussian Splatting) can thiệp: bổ sung thêm một tín
    hiệu đo trực tiếp \textbf{cấu trúc tần số cục bộ} của ảnh, để quyết định "đúng chỗ, đúng lúc, đúng trục" hơn.
\end{itemize}
\vspace{2pt}
\footnotesize\textbf{Lộ trình bài giảng:} (1) Biểu diễn Gaussian \& pipeline render $\to$ (2) Loss, gradient, ADC nguyên
bản của 3DGS $\to$ (3) SADGS: xây dựng tín hiệu tần số $\eta$ từ đầu, cơ chế densify/prune mới $\to$ (4) Nhìn lại một
cách trung thực: cái gì đang thật sự chạy, so sánh với bản gốc.
\end{frame}

% --- Slide: 3D Gaussian là gì ---
\begin{frame}{3D Gaussian: đơn vị biểu diễn cơ bản}
\scriptsize
Toàn bộ cảnh được biểu diễn bằng một \textbf{tập hợp các Gaussian 3D} --- mỗi Gaussian là một hàm mật độ liên tục
trong không gian, tâm tại $\mu\in\mathbb{R}^3$, "độ loang" theo ma trận hiệp phương sai $\Sigma\in\mathbb{R}^{3\times3}$:
\[
\boxed{\;G(x)=\exp\!\Big(-\tfrac12(x-\mu)^\top\Sigma^{-1}(x-\mu)\Big)\;}
\]
Có thể hình dung mỗi Gaussian như một \textbf{"cục mây mờ" hình ellipsoid}: đậm đặc nhất ở tâm $\mu$, mờ dần ra biên
theo hình dạng do $\Sigma$ quyết định. Mỗi Gaussian mang theo $4$ nhóm tham số \textbf{học được bằng gradient descent}
(tổng $59$ số thực: xyz 3 $+$ scaling 3 $+$ rotation 4 $+$ opacity 1 $+$ \texttt{f\_dc} 3 $+$ \texttt{f\_rest} 45):
\begin{itemize}\itemsep1pt
    \item $\mu$ --- vị trí tâm trong không gian 3D (3 số).
    \item $\Sigma$ --- hình dạng và hướng của ellipsoid, không lưu trực tiếp mà tham số hoá qua scale $s\in\mathbb{R}^3$
    và quaternion (hướng quay) $q\in\mathbb{R}^4$ --- lý do tại sao xem ở slide tiếp theo.
    \item $\alpha\in(0,1)$ --- độ mờ đục (opacity): Gaussian "đặc" hay "trong suốt", lưu dưới dạng logit rồi đi qua sigmoid để luôn nằm trong khoảng hợp lệ.
    \item $c$ --- màu sắc, nhưng \textbf{không phải một màu RGB cố định} mà biểu diễn bằng hệ số Spherical Harmonics
    (Slide sau) để màu có thể \emph{đổi theo góc nhìn} --- mô phỏng hiệu ứng phản chiếu, ánh kim, specular thực tế.
\end{itemize}
\begin{center}
\includegraphics[width=0.4\linewidth]{ch2_gaussian_bump.png}
\end{center}
\end{frame}

% --- Slide: Tham số hoá hiệp phương sai ---
\begin{frame}{Tham số hoá hiệp phương sai: luôn xác định dương}
\scriptsize
Về mặt toán học, $\Sigma$ phải \textbf{đối xứng và xác định dương} mới là một ma trận hiệp phương sai hợp lệ (nếu
không, $G(x)$ ở trên có thể phát tán vô nghĩa). Nếu tối ưu trực tiếp $6$ số độc lập của $\Sigma$ bằng gradient descent,
rất dễ trong quá trình huấn luyện đi lệch khỏi vùng hợp lệ. Giải pháp của 3DGS là \textbf{phân rã} $\Sigma$ thành tích
của một ma trận quay và một ma trận scale:
\[
\boxed{\;\Sigma=RSS^\top R^\top\;},\qquad S=\mathrm{diag}(s_x,s_y,s_z),\qquad R=R\!\left(\frac{q}{\lVert q\rVert}\right)
\]
\begin{itemize}\itemsep1pt
    \item $S$: ma trận đường chéo với $s=\exp(\texttt{\_scaling})$ --- vì hàm mũ luôn dương, $s$ \textbf{tự động dương}
    với mọi giá trị thực của \texttt{\_scaling}, không cần ràng buộc (constraint) nào thêm khi tối ưu.
    \item $R$: ma trận quay $3\times3$ dựng từ quaternion $q$ đã \textbf{chuẩn hoá về độ dài 1} --- luôn là ma trận trực giao hợp lệ.
    \item Nhờ vậy $\Sigma=RSS^\top R^\top$ \textbf{tự động} đối xứng, xác định dương với \emph{bất kỳ} $s,q$ nào --- ta có
    thể tối ưu $s,q$ tự do bằng gradient descent mà không sợ vi phạm ràng buộc toán học.
    \item Về mặt hình học: ellipsoid của Gaussian có \textbf{bán trục theo mỗi hướng chính $= s_j$} (không phải trị riêng
    của $\Sigma$ --- trị riêng thực ra là $s_j^2$, vì $\Sigma$ là bình phương của $S$ sau khi quay).
\end{itemize}
\end{frame}

% --- Slide: Scale khởi tạo ---
\begin{frame}{Khởi tạo (1/2): scale ban đầu theo mật độ điểm SfM}
\scriptsize
Trước khi train, mỗi Gaussian được đặt \textbf{đúng tại vị trí} một điểm 3D thưa mà SfM/COLMAP đã tái dựng được từ các
ảnh input (\emph{sparse point cloud}). Nhưng vị trí thôi chưa đủ --- cần một \textbf{kích thước ban đầu hợp lý} cho mỗi
Gaussian: quá nhỏ thì để hở nhiều khoảng trống giữa các điểm, quá to thì các Gaussian chồng lấn nhau ngay từ đầu, làm
nhiễu tín hiệu học. Giải pháp: lấy khoảng cách tới các điểm lân cận làm thước đo mật độ cục bộ:
\[
\tilde s = \log\sqrt{\overline{d^2_{\mathrm{knn3}}}}
\]
tức là: với mỗi điểm, tính khoảng cách bình phương tới $3$ điểm gần nhất (knn3), lấy trung bình, rồi căn bậc hai --- ra
một ước lượng "khoảng cách điển hình giữa các điểm lân cận" tại vị trí đó. Lấy \texttt{log} vì scale được lưu và tối
ưu trong không gian log (nhất quán với $s=\exp(\texttt{\_scaling})$ ở slide trước). Kết quả: vùng point cloud dày đặc
(nhiều chi tiết hình học) sẽ có Gaussian khởi tạo \textbf{nhỏ}, vùng thưa sẽ có Gaussian khởi tạo \textbf{to} hơn ---
một cách tự nhiên để phủ kín cảnh mà không cần biết trước hình dạng vật thể.
\begin{center}
\includegraphics[width=0.42\linewidth]{ch1_scale_init.png}
\end{center}
\end{frame}

% --- Slide: Scene extent ---
\begin{frame}{Khởi tạo (2/2): scene extent --- thước đo "to/nhỏ" dùng xuyên suốt}
\scriptsize
Nhiều quyết định sau này trong ADC (Adaptive Density Control) cần biết "Gaussian này to hay nhỏ \textbf{so với quy mô
của cả cảnh}" --- ví dụ ngưỡng \texttt{percent\_dense}$\times$\texttt{extent} để phân biệt clone/split (sẽ nói ở phần
sau). Đại lượng dùng làm đơn vị đo lường này là \textbf{scene extent}, tính \textbf{một lần duy nhất} lúc khởi tạo:
\[
\boxed{\;\mathrm{extent}=1.1\cdot\max_v\lVert c_v-\bar c\rVert\;}
\]
với $c_v$ là vị trí của camera train thứ $v$, $\bar c$ là tâm trung bình của tất cả camera train. Nói cách khác, đây là
\textbf{bán kính} (nhân thêm hệ số an toàn $1.1$) của một khối cầu tưởng tượng bao quanh toàn bộ vị trí các camera ---
\textbf{không phải} đường chéo của hộp bao (bounding box) quanh point cloud như đôi khi bị hiểu nhầm. Vì cảnh thường
được quay quanh một đối tượng trung tâm, "bán kính vòng camera" là một ước lượng hợp lý cho kích thước tổng thể của cảnh.
\begin{center}
\includegraphics[width=0.42\linewidth]{ch1_extent.png}
\end{center}
\end{frame}

% --- Slide: Màu qua Spherical Harmonics ---
\begin{frame}{Màu phụ thuộc góc nhìn: hệ số Spherical Harmonics}
\scriptsize
Vật liệu thực tế (kim loại, nước, sứ bóng...) phản chiếu ánh sáng khác nhau tuỳ góc nhìn --- một màu RGB cố định không
mô phỏng được hiệu ứng này. 3DGS giải quyết bằng cách biểu diễn màu như một \textbf{hàm của hướng nhìn} $d$, khai triển
theo cơ sở \textbf{Spherical Harmonics} (SH) --- một họ hàm trực giao trên mặt cầu, tương tự chuỗi Fourier nhưng cho dữ
liệu định nghĩa trên hình cầu thay vì trên một đường thẳng:
\[
c(d)=\sum_{\ell=0}^{D}\sum_{m=-\ell}^{\ell} k_{\ell m}\,Y_{\ell}^{m}(d)
\]
mỗi Gaussian lưu các \textbf{hệ số} $k_{\ell m}$ (học được), không lưu trực tiếp màu.
\begin{itemize}\itemsep1pt
    \item Hệ số bậc $0$ ($\ell=0$, gọi là \textbf{hệ số DC}) chính là \textbf{màu trung bình theo mọi góc nhìn}: khởi tạo
    $k_{00}=\mathrm{RGB2SH}(c)=(c-0.5)/C_0$ với hằng số chuẩn hoá $C_0=\tfrac{1}{2\sqrt\pi}\approx0.282$.
    \item Các hệ số bậc cao hơn ($\ell\ge1$) mã hoá \textbf{độ thay đổi màu theo hướng nhìn} --- càng nhiều bậc, hiệu ứng
    phản chiếu/specular càng chi tiết, nhưng cũng càng nhiều tham số phải học.
    \item Tổng số hệ số cần học nếu dùng tới bậc $D$: $(D+1)^2$ hệ số cho mỗi kênh màu; với $D=3$ (giá trị dùng phổ biến)
    là $16$ hệ số/kênh $\times3$ kênh $=48$ số ($3$ số DC $+$ $45$ số bậc cao, tên biến \texttt{f\_rest}).
\end{itemize}
\begin{center}
\includegraphics[width=0.30\linewidth]{ch1_rgb2sh.png}
\end{center}
\end{frame}

% --- Slide: Lịch tăng bậc SH ---
\begin{frame}{Lịch tăng bậc SH: vì sao không học bậc cao ngay từ đầu?}
\scriptsize
Học ngay $48$ hệ số SH từ những iteration đầu tiên, khi vị trí/hình dạng Gaussian còn rất chưa ổn định, dễ làm mô hình
"khớp" (overfit) vào những biến thiên màu do \textbf{sai vị trí hình học} gây ra, chứ không phải phản chiếu thật ---
gây nhiễu và bất ổn cho quá trình tối ưu. Giải pháp: \textbf{tăng bậc SH dần dần}, để hình học ổn định trước, chi tiết
phản chiếu học sau. Trong code (\texttt{train.py:216--217}):
\[
\texttt{if iteration \% 1 == 0: gaussians.oneupSHdegree()}
\]
tức bậc \textbf{tăng thêm 1 sau mỗi iteration}, cho tới khi chạm bậc tối đa $D_{\max}=3$ --- nghĩa là chỉ sau \textbf{3
iteration đầu tiên} của cả quá trình train ($30\,000$ iteration), toàn bộ Gaussian đã dùng bậc SH tối đa. Đây là một
điểm dễ hiểu lầm: nhiều tài liệu khác về 3DGS mô tả lịch tăng bậc là "mỗi vài nghìn iteration tăng 1 bậc" --- con số đó
\textbf{không khớp} với đoạn code thực tế của chính pipeline này.
\begin{center}
\includegraphics[width=0.42\linewidth]{ch2_sh_degree_schedule.png}
\end{center}
\end{frame}
